\title{Parameter-Efficient Orthogonal Finetuning via Factorization}

\begin{abstract}
While large foundation models have become widespread, retraining them from the ground up is extremely costly. As a result, finding efficient methods to tailor these robust models to specific downstream tasks has gained critical significance. This work explores Orthogonal Finetuning (OFT), a theoretically-grounded approach for task adaptation. Although OFT has shown strong transferability, its reliance on orthogonal matrices with substantial dimensionality still leads to a high number of parameters to be trained. To mitigate this, we first analyze OFT through the lens of information transfer and derive several fundamental criteria for improving parameter efficiency. Drawing inspiration from the Cooley-Tukey fast Fourier transform’s effectiveness in optimizing information flow, we introduce a compact orthogonal parameterization using butterfly structures. Integrating this parameterization into OFT yields a new, more parameter-efficient approach, named Orthogonal Butterfly (BOFT). BOFT not only encompasses OFT as a specific instance but also offers a broader framework for orthogonal finetuning. We validate BOFT through extensive experiments, adapting a range of large vision transformers, large language models, and text-to-image diffusion models to numerous vision and language downstream tasks.
\end{abstract}

\section{Introduction}

State-of-the-art models like ChatGPT~\cite{brown2020language,chen2021evaluating} and Stable Diffusion~\cite{rombach2022high} showcase the impressive generalization capabilities inherent in large foundation models. This swift advancement in model performance is closely linked to a substantial escalation in parameter counts.

(\emph{e.g.}, GPT-3 comprises approximately 175B parameters). Consequently, building a foundation model from the ground up has become increasingly prohibitive for researchers. As the field advances, it is crucial to develop techniques for modifying these models without resorting to full retraining. In essence, the ability to efficiently tailor foundation models for specific downstream applications is essential. There are three predominant strategies for such efficient task adaptation: \emph{model finetuning}~\cite{hulora2022,zaken2022bitfit,guo2020parameter,raffel2020exploring,qiu2023controlling,zhang2022adaptive,chavan2023one}, where only a portion of the model’s parameters are updated while the original architecture remains intact; \emph{adapter tuning}~\cite{rebuffi2017learning,houlsby2019parameter,pfeiffer2020adapterfusion,lin2020exploring,he2021towards}, which involves augmenting the model with additional trainable components; and \emph{prompt tuning}~\cite{li2021prefix,lester2021power}, in which supplementary trainable prefix tokens are appended to the model’s input. Among these techniques, model finetuning stands out for its simplicity and effectiveness, as well as its advantage of introducing no additional latency during inference.

The main concept underlying model finetuning is to maintain a high degree of similarity between the pretrained and finetuned models according to specific metrics, thereby retaining the knowledge acquired during pretraining. In practice, most contemporary finetuning approaches utilize a low learning rate as a heuristic to constrain the deviation between the initial and final models. Recognizing the inherent structure within weight matrices, a more systematic strategy aims to conserve the relational characteristics—specifically, the pairwise angular relationships among neurons. This perspective motivates the development of Orthogonal Finetuning~(OFT)~\cite{qiu2023controlling}, a framework in which each neuron within a layer is transformed using the same orthogonal matrix, guaranteeing that the angular relationships between neurons are mathematically maintained throughout finetuning. While OFT has yielded notable improvements in both generalizability and convergence, particularly when applied to the finetuning of text-to-image diffusion models~\cite{qiu2023controlling}, a significant limitation is the substantial parameter count arising from the high dimensionality inherent in orthogonal matrices. To mitigate this, OFT incorporates a block-diagonal design that lowers the number of learnable parameters. Nonetheless, this increase in parameter efficiency comes with certain drawbacks: the block-diagonal orthogonal matrices exhibit a fixed pattern of sparsity, and transformations occur independently within each block. Such structural sparsity can introduce unwelcome inductive biases (for instance, limiting the expressive capacity of the orthogonal transformations and precluding the approximation of traditional linear mappings). Additionally, the question of how to select an effective sparsity configuration remains unresolved.

The central challenge is to construct a dense orthogonal matrix in a way that remains efficient in terms of the number of parameters. Although it initially appears that a dense $d$-dimensional orthogonal matrix necessitates $\mathcal{O}(d^2)$ parameters, our method departs from convention by assembling such a matrix through a composition of several factorized sparse matrices. The foundational observation behind our approach is that the parameter count can be minimized by increasing computational effort rather than memory usage. Because the orthogonal matrix is expressed as a product of sparse factors, each training step requires successive multiplication by these sparse matrices. To better conceptualize the matrix factorization process, we view constructing a dense orthogonal matrix as analogous to transmitting information across a network with a grid topology. In this context, building up a dense orthogonal matrix from products of sparse matrices corresponds to relaying information throughout the grid-structured graph. Adopting this information transmission lens offers a systematic way to explore diverse strategies for sparse matrix factorization, allowing us to maintain a manageable number of trainable parameters while ensuring the resulting matrices can achieve the desired density and expressive capacity.

In pursuit of parameter efficiency within our information transmission scheme, we take cues from the butterfly structures employed in the Cooley-Tukey fast Fourier transform algorithm~\cite{cooley1965algorithm}, known for their effectiveness in facilitating information flow between distinct nodes~\cite{klasing1994broadcasting}. The butterfly graph central to the Cooley-Tukey algorithm naturally leads to a type of sparse matrix decomposition, termed butterfly factorization. Through this factorization, a dense matrix of dimension $d$ can be represented as a product of $\mathcal{O}(\log d)$ sparse matrices, where each matrix contains only $\mathcal{O}(d)$ nonzero elements. By enforcing orthogonality on each of these sparse factors, we derive an efficient orthogonal parameterization requiring merely $\mathcal{O}(d\log d)$ parameters, representing a drastic reduction relative to the original dense case. Building on this compact orthogonal parameterization, we introduce Orthogonal Butterfly~(BOFT), a new parameter-efficient finetuning strategy. BOFT encompasses OFT as a particular instance, and more broadly establishes a unified orthogonal finetuning paradigm. Both the block-diagonal and butterfly matrix constructions are characterized by their sparsity, which injects sparsity into orthogonal matrices, thereby diminishing the number of parameters necessitating training. A comparison with LoRA~\cite{hulora2022}, which leverages low-rank parameterizations, is instructive: sparse matrices and low-rank matrices constitute two prominent categories of structured matrices~\cite{candes2011robust} that are widely utilized to achieve parameter-efficient representations.

In contrast to the block-diagonal design employed by OFT, which negotiates between expressiveness and structural regularity, BOFT leverages the butterfly configuration to offer a more continuous transition between full orthogonal group matrices (offering expressiveness) and the identity matrix (providing regularity). This construction facilitates the identification of a more constrained hypothesis space within the orthogonal group that is particularly useful for subsequent tasks. Given the butterfly structure’s prevalent role in various efficient linear transformations—including the discrete Fourier transform and discrete cosine transform—we propose that our structural approach to parameter efficiency imparts an implicit inductive bias that can improve generalization and reduce overfitting. This perspective is motivated by the butterfly structure's capability to replicate many foundational linear transformations, a property absent in OFT's block-diagonal form. Below, we summarize our main contributions:

\item We investigate parameter efficiency in orthogonal finetuning by introducing a new information transmission perspective. Within this framework, the challenge of designing a parameter-efficient dense orthogonal matrix is reimagined as an information flow optimization problem over a graph with a grid topology.

\item Drawing inspiration from the butterfly factorization used in the Cooley-Tukey algorithm, we introduce Orthogonal Butterfly—a novel orthogonal finetuning technique with high parameter efficiency—constructed in accordance with our information transmission framework.

\item We present theoretical justifications for how BOFT achieves substantial reduction in the count of trainable parameters, yet maintains strong expressive power and training robustness. The matrix decomposition underpinning BOFT naturally facilitates a compelling approach to weight interpolation (see Figure~\ref{fig:boft_interpolation}).

\item For the first time, we extend orthogonal finetuning to a diverse array of tasks outside of controllable text-to-image generation~\cite{qiu2023controlling}, highlighting its promise as a general finetuning technique. In particular, we demonstrate the effectiveness of BOFT across multiple adaptation tasks, spanning domains such as computer vision and natural language processing. Our results show that BOFT surpasses existing leading approaches by a significant margin, confirming its outstanding parameter efficiency and generalization performance.

\section{Related Work}

\textbf{Parameter-efficient finetuning (PEFT).} As the scale and capability of foundation models (\emph{e.g.}, \cite{brown2020language,radford2021learning,rombach2022high,kirillov2023segment}) continue to grow, there has been substantial interest in strategies that enable efficient finetuning for specific applications while minimizing the number of parameters modified~\cite{treviso2023efficient,ding2023parameter,lialin2023scaling}. Numerous PEFT approaches have been introduced~\cite{houlsby2019parameter,aghajanyan2020intrinsic,hulora2022,edalati2022krona,wang2022adamix,gheini2021cross,zaken2022bitfit,guo2020parameter,sung2021training,ansell2022composable,lester2021power,li2021prefix,vu2022spot,he2021towards,mao2021unipelt,karimi2021compacter,liu2022few,sung2022lst,chen2022parameter,jia2022visual,chen2022adaptformer,zhang2022neural,jie2023fact,lian2022scaling,luo2023towards}, with reparameterization techniques~\cite{aghajanyan2020intrinsic,hulora2022,edalati2022krona} being especially pertinent to the present study. LoRA~\cite{hulora2022} modifies pretrained weight matrices by incorporating the product of two low-rank matrices, leading to competitive outcomes on natural language processing tasks. While LoRA fixes the rank for these added matrices, subsequent research \cite{zhang2022adaptive,valipour2022dylora,zhang2023increlora} has explored adaptively varying the rank across different layers to optimize parameter allocation. The appeal of low-rank reparameterization approaches stems largely from their ease of implementation, which has contributed to their widespread adoption \cite{dettmers2023qlora,zi2023delta,chavan2023one}. Motivated by studies on hyperspherical energy and its connection to generalization~\cite{liu2018learning,liu2021orthogonal}, \cite{qiu2023controlling} put forward orthogonal finetuning, which introduces an orthogonal transformation within each layer as a means of finetuning text-to-image diffusion models. This technique, which optimizes an orthogonal matrix for intra-layer neuron transformation, not only offers improved generalization but also enhances training stability compared to LoRA. Nevertheless, orthogonal finetuning (OFT) generally results in a higher parameter count than LoRA. As such, enhancing the parameter efficiency of OFT is a pertinent objective. In addition, it remains to be determined whether OFT can be generalized beyond text-to-image diffusion control tasks~\cite{qiu2023controlling}. BOFT addresses OFT’s parameter overhead by employing butterfly factorization. This advancement enables us to extend orthogonal finetuning to a broader array of adaptation scenarios.

\textbf{Butterfly structures}. The radix-2 Cooley-Tukey algorithm achieves the decomposition of an $N$-point discrete Fourier transform by recursively breaking it down into two discrete Fourier transforms of size $\frac{N}{2}$. This recursive division gives rise to a distinctive butterfly structure, which is representable as a product of several sparse matrices—commonly referred to as a butterfly matrix. Prior work has leveraged butterfly matrices to parameterize orthogonal matrices for circumventing the need for pivoting in Gaussian elimination and for improving computational efficiency~\cite{parker1995random}, as well as for promoting stable recurrent neural network training~\cite{jing2017tunable} and kernel approximation~\cite{munkhoeva2018quadrature}. Methods in~\cite{dao2019learning,dao2019kaleidoscope} develop efficient linear mappings by adopting butterfly parameterizations, while~\cite{chen2022pixelated,dao2022monarch} employ butterfly matrices to realize efficient neural network training procedures. Butterfly structures have also demonstrated benefits in applications such as accelerated matrix-vector products~\cite{michielssen1996multilevel,o2010algorithm}, data-sparse matrix representations~\cite{li2015butterfly}, and in the context of network transmission~\cite{klasing1994broadcasting,li2003linear,rajasingh2016transmission}. Distinct from previous studies, our approach centers on leveraging butterfly structures to improve the parameter efficiency of OFT.

\section{An Information Transmission View on Orthogonal Finetuning}

We begin by reviewing key concepts underlying OFT. In order to adapt pretrained weights, OFT reformulates the weight matrix as the product of a trainable orthogonal matrix and the fixed pretrained weights. Unlike LoRA, which incorporates an additive low-rank adjustment to the weights, OFT applies a multiplicative orthogonal matrix for weight adaptation. LoRA’s parameter efficiency arises from the use of low-rank updates, whereas classic OFT employs a block-diagonal orthogonal structure~\cite{qiu2023controlling}; efficiency improves as the blocks become smaller. Figure~\ref{fig:comp_lora} provides a side-by-side comparison. The rationale for adopting orthogonal transformations in OFT is to retain the angular relationships between neurons~\cite{liu2018learning,liu2021orthogonal,qiu2023controlling}, which in turn helps preserve semantic information acquired during pretraining. 

Specifically, OFT learns an orthogonal matrix $\bm{R}\in\mathbb{R}^{d\times d}$ in conjunction with a pretrained linear layer $\bm{W}^0\in\mathbb{R}^{d\times n}$, replacing the forward computation $\bm{z}=(\bm{W}^0)^\top\bm{x}$ with $\bm{z}=(\bm{R}\bm{W}^0)^\top\bm{x}$, where the input vector $\bm{x}\in\mathbb{R}^{d}$ maps to output $\bm{z}\in\mathbb{R}^{n}$. For zero initialization, $\bm{R}$ is set to the identity matrix at the outset. To maintain orthogonality of $\bm{R}$ during training, we utilize Cayley parameterization as in \cite{qiu2023controlling,liu2021orthogonal}: $\bm{R}=(\bm{I}+\bm{Q})(\bm{I}-\bm{Q})^{-1}$, where $\bm{Q}$ is a skew-symmetric matrix satisfying $\bm{Q}=-\bm{Q}^\top$. To further enhance parameter efficiency, a block-diagonal form is adopted, parameterizing $\bm{R}$ as $\text{diag}(\bm{R}_1,\bm{R}_2,\cdots,\bm{R}_r)$ with each $\bm{R}_i\in\mathbb{R}^{b\times b}$ representing a compact orthogonal matrix and $br=d$. However, the block-diagonal approach imposes an implicit constraint: each neuron's vector (a column of $\bm{W}^0$) is partitioned into $r$ disjoint segments, with distinct orthogonal matrices applied independently to each group. While the block-diagonal design proves effective in experiments, it is theoretically unmotivated to segment a neuron's vector space solely by index, highlighting the appeal of dense orthogonal matrices. This leads to the key question: \emph{Is it possible to obtain a dense orthogonal matrix while retaining parameter efficiency?}

In response to this question, we suggest representing a dense orthogonal matrix as the product of several sparse orthogonal matrices. To offer both a cohesive and intuitive framework for analyzing the sparsity structures involved in orthogonal matrix factorization, we reinterpret the construction of a dense orthogonal matrix in OFT as analogous to an information transmission process. Concretely, constructing a dense matrix $\bm{R}\in\mathbb{R}^{d\times d}$ from the product of $m$ square matrices, $\thickmuskip=2mu \medmuskip=2mu \bm{R}=\bm{B}_m\bm{B}_{m-1}\cdots\bm{B}_1$, can be visualized as facilitating information transfer across a lattice with $d\times (m+1)$ nodes, as depicted in Figure~\ref{fig:info_trans}. This viewpoint arises from recognizing that a $d$-dimensional dense square matrix effectively describes full connectivity from a set of $d$ source nodes to another set of $d$ target nodes. In this representation, an entry $R_{ij}$ being zero in $\bm{R}$ indicates the absence of a transmission pathway from the $j$-th source node to the $i$-th target node, whereas a nonzero $R_{ij}$ suggests an active communication link. Accordingly, decomposing $\bm{R}$ into the product $\bm{B}_m\bm{B}_{m-1}\cdots\bm{B}_1$ can be viewed as a sequence of staged information transfers determined by the connectivity graphs corresponding to each $\bm{B}_i$. Information propagation follows the structure imposed by $\bm{B}_1$ at the start and concludes via $\bm{B}_n$. 

A concrete illustration is given in Figure~\ref{fig:info_trans} through the factorization $\bm{R}=\bm{B}_5\bm{B}_4\bm{B}_3\bm{B}_2\bm{B}_1$, with the sparsity structures and corresponding graph displayed. This induced graph emerges from expanding the series of matrix multiplications. In such a graph, each matrix $\bm{B}_i$ functions as a connectivity blueprint linking nodes between consecutive layers, from the $i$-th level to the $(i+1)$-th. More precisely, an entry $(j_1, j_2)$ in $\bm{B}_i$ signifies whether there is a directed edge from the $j_2$-th node of the $i$-th level to the $j_1$-th node of the $(i+1)$-th level; a zero indicates the absence of such a connection. 

For $\bm{R}=\bm{B}_5\bm{B}_4\bm{B}_3\bm{B}_2\bm{B}_1$ to yield a dense matrix, each initial node must be capable of sending information to all nodes in the $6$-th (final) layer. In contrast, for $\bm{R}=\bm{B}_4\bm{B}_3\bm{B}_2\bm{B}_1$, which models transmissions from first to fifth levels, it may occur that, for instance, node $1$ at the source cannot communicate with node $3$ at the destination, corresponding to a zero at position $(3,1)$ in the sparsity pattern of the composite matrix. This same mapping between graph reachability and sparsity structure persists in shorter products, such as $\bm{R}=\bm{B}_3\bm{B}_2\bm{B}_1$ and $\bm{R}=\bm{B}_2\bm{B}_1$. 

More generally, to ensure $\bm{R}\in\mathbb{R}^{d\times d}$ is dense, the collection of $m$ factor matrices $\bm{B}_m,\cdots,\bm{B}_1$ must collectively define a network of directed edges within a $d\times (m+1)$ grid. Each directed edge only connects two nodes in adjacent layers (that is, between neighboring columns), and this edge arrangement must be such that every node in the initial layer has a transmission route to every node in the $(m+1)$-th layer.

Adopting the perspective of information transmission provides valuable insights into the sparsity patterns that arise within the factorization matrices used in OFT. As shown in Figure~\ref{fig:blk_diag}, $\bm{R}$ in the original OFT exhibits a distinct block-diagonal form. While this structure helps to limit the count of trainable parameters, it is inherently incapable of yielding a fully dense matrix $\bm{R}$. Our objective is to generate a dense orthogonal matrix by composing $m$ sparse orthogonal matrices, achieving this with the minimal necessary number of effective trainable parameters.

From the information transmission standpoint, two primary criteria emerge: \emph{(i)} \textbf{dense connectivity}, meaning that each node at the initial layer is connected—by at least one pathway—to every node in the final layer; and \emph{(ii)} \textbf{minimum free edges}, meaning the aggregate edge count should be minimized, subject to maintaining orthogonality. Notably, the orthogonality constraint imposes specific structural requirements on the edges between consecutive layers. For instance, to guarantee that each matrix $\bm{B}_i$ is full-rank—a prerequisite for orthogonality—there must be $d$ edges forming a bijection between all nodes of level $i$ and those of level $(i+1)$. This necessitates a minimum of $d$ inter-layer edges (see, for example, Figure~\ref{fig:info_trans}, where $d=4$). Importantly, these $d$ edges are dictated by orthogonality, are not adjustable during training, and thus are excluded from the count of trainable edges (e.g., the nonzero entries in a $d\times d$ orthogonal matrix are fixed at $\pm 1$). 

Orthogonal matrices require fewer free parameters than unrestricted ones, meaning the orthogonality constraint introduces extra dependencies between the possible edges. As a case in point, setting a zero in the $(1,1)$ entry of a $2\times2$ orthogonal matrix automatically dictates a zero at the $(2,2)$ entry—so the omission of one edge can entail the absence of another. Hence, depending on the feasible configuration of connections, orthogonality may lead to the inclusion or exclusion of certain edges. By examining the locations of nonzero entries in the resulting orthogonal matrix, the information transmission approach offers particular clarity for OFT scenarios, where our concern centers only on the placement of trainable, nonzero entries in $\bm{R}$, with their specific magnitudes being of lesser relevance.

A straightforward approach to densely connect two levels requires $\mathcal{O}(d^2)$ edges (that is, it corresponds to a single dense orthogonal matrix), which equates to $d^2 - d$ edges; for instance, when $d=4$, this totals 12 edges. In Figure~\ref{fig:info_trans}, we illustrate an example of a practical matrix factorization, which uses just $10$ edges—fewer than what a full dense orthogonal matrix demands. By adopting this graphical perspective, we are able to analyze the parameter-efficiency properties of OFT, and this viewpoint facilitates the design of viable matrix factorizations. Our methodology draws on an intriguing topology used in the Cooley-Tukey algorithm known as butterfly graphs~\cite{cooley1965algorithm}. Butterfly graphs are capable of establishing efficient dense connections between $d$ input and $d$ output nodes, requiring only $\mathcal{O}(d\log d)$ edges. For context, while the structure in Figure~\ref{fig:info_trans} needs $10$ edges for complete connectivity, a butterfly network would need just $8$. In the following section, we discuss how leveraging the butterfly structure can lead to further improvements in parameter efficiency.

\section{Orthogonal Parameterization by Butterfly Factorization}

The butterfly structure was initially introduced in the Cooley-Tukey algorithm for executing fast Fourier transforms. In the context of Fourier transforms, alterations that are localized in the frequency domain may induce widespread effects across the spatial domain. This relationship parallels the information transmission scenario under discussion, where each node at the initial stage is capable of delivering information to any node at the terminal stage. Moreover, the butterfly structure has gained widespread adoption as an efficient topology for computer networks~\cite{solihin2015fundamentals,li2011linear} to facilitate information exchange. Given $k\geq 2$ as a power of $2$, we first formalize the \emph{butterfly factor} $\bm{B}^F(k)$ as

\begin{equation}\label{eq:but_factor}
\footnotesize
\begin{aligned}
    \bm{B}^F(k)=\begin{bmatrix}
    \text{diag}(\bm{d}_1) & \text{diag}(\bm{d}_2) \\
    \text{diag}(\bm{d}_3) & \text{diag}(\bm{d}_4)
    \end{bmatrix}\in\mathbb{R}^{k\times k},
\end{aligned}
\end{equation}

where each $\bm{d}_i\in\mathbb{R}^{\frac{k}{2}}$ for all $i$ represent vectors. For $d=2^N$, we introduce the $d$-dimensional \emph{butterfly matrix} $\bm{B}(d)\in\mathbb{R}^{d\times d}$, which is constructed recursively as follows:

\begin{equation}
\footnotesize
\begin{aligned}
    \!\!\bm{B}(d)=\tilde{\bm{B}}(d,d)\!\cdot\!\begin{bmatrix}
    \bm{B}_1(\frac{d}{2})\!\!\!\!\! & \bm{0} \\
    \bm{0}\!\!\!\!\! &\bm{B}_2(\frac{d}{2})
    \end{bmatrix}= \tilde{\bm{B}}(d,d)\tilde{\bm{B}}(d,\frac{d}{2})\cdots\tilde{\bm{B}}(d,2),
\end{aligned}
\end{equation}

where $\bm{B}_1(\frac{d}{2})$ and $\bm{B}_2(\frac{d}{2})$ denote two butterfly matrices of dimension $\frac{d}{2}$. We subsequently introduce the \emph{butterfly component}, denoted as $\thickmuskip=2mu \medmuskip=2mu \tilde{\bm{B}}(d,k)=\text{diag}(\bm{B}^F_1(k),\cdots,\bm{B}^F_{d/k}(k))$, representing a block-diagonal matrix of shape $d\times d$ with individual blocks of size $k$. Here, each $\bm{B}^F_i(k)$ is a butterfly factor, as defined in Equation~\ref{eq:but_factor}. This construction lays the foundation for utilizing butterfly matrices as a means to parametrize orthogonal matrices. For this purpose, it is necessary that every factor $\tilde{\bm{B}}(d,k)$ in the butterfly matrix $\bm{B}(d)$ is orthogonal. 

Focusing on the block-diagonal structure $\tilde{\bm{B}}(d,2)$ where the block size is $2$, orthogonality is straightforwardly enforced by parametrizing each block using the Cayley transform (or equivalently, a $2$-dimensional rotation), as is also done in \cite{liu2021orthogonal,qiu2023controlling}. Importantly, the arrangement of non-zero elements in any butterfly component is merely a permuted version of that in $\tilde{\bm{B}}(d,2)$, allowing for each butterfly component to be parameterized as an orthogonal matrix in a similar manner. This approach efficiently builds an orthogonal matrix parameterization using numerous $2\times 2$ orthogonal submatrices. 

Extending the butterfly matrix approach following \cite{chen2022pixelated}, we define a block butterfly component $\tilde{\bm{B}}^b(d,k)$, in which each element of $\bm{d}_i$ is a $b\times b$ matrix. To enforce orthogonality for the block butterfly component $\tilde{\bm{B}}^b(d,2)$, we parametrize every $2b\times 2b$ block matrix as orthogonal. For butterfly components $\tilde{\bm{B}}^b(d,k)$ with $k>2$, their non-zero patterns are block-wise permutations of those in $\tilde{\bm{B}}^b(d,2)$, so these too can be made orthogonal using the same technique. Assembling these elements, the BOFT forward pass is given by

\begin{equation}
    \footnotesize
    \begin{aligned}\nonumber
    \bm{z}=\big(\bm{R}(m,b)\cdot\bm{W}^0\big)^\top\bm{x},~~~~\text{s.t.}~\bigg{\{}\bm{R}(m,b)=\prod_{i=1}^m \tilde{\bm{B}}^b_{(d,i)}~~\&~~\underbrace{\big(\tilde{\bm{B}}^b_{(d,j)}\big)^\top\tilde{\bm{B}}^b_{(d,j)}=\tilde{\bm{B}}^b_{(d,j)}\big(\tilde{\bm{B}}^b_{(d,j)}\big)^\top\!=\bm{I}_d}_{\forall j\in [1, m]}\bigg{\}},
    \end{aligned}
\end{equation}

Here, for notational brevity, we write $\tilde{\bm{B}}^b(d,2^{m - i+1})$ as $\tilde{\bm{B}}^b_{(d,i)}$, and $\bm{I}_d$ denotes the $d$-dimensional identity matrix. The orthogonal matrix $\bm{R}(m,b)\in\mathbb{R}^{d\times d}$ is structured as a sequence of orthogonal butterfly components multiplied together. We refer to BOFT with $\bm{R}(m,\frac{b}{2})$ simply as BOFT($m$,$b$) for all $b \geq 2$. Notably, when $m=1$, BOFT($1$,$b$) collapses to the block-diagonal OFT~\cite{qiu2023controlling} having block size $b$. Likewise, when $b=d$, BOFT($1$,$d$) becomes the standard OFT~\cite{qiu2023controlling} that employs an unrestricted full orthogonal matrix. Employing the block butterfly matrix $\bm{B}^{\frac{b}{2}}(d)$ as $\bm{R}$ with BOFT($\log \frac{2d}{b}$,$b$) results in a dense orthogonal matrix $\bm{R}$. More generally, BOFT($m$,$b$) introduces a total of $\frac{1}{2}(b-1)dm$ effective trainable parameters when finetuning a linear mapping of size $d\times n$. Selecting the butterfly construction with $m=\log d$ and $b=2$ reduces the parameter count for BOFT to $\mathcal{O}(d\log d)$. By comparison, the traditional OFT utilizing a fully dense orthogonal matrix requires $\mathcal{O}(d^2)$ parameters, while the block-diagonal OFT with $r$ blocks incurs $\mathcal{O}(bd)$. Consequently, obtaining a dense orthogonal matrix via the original OFT demands $b=d$ for block size, whereas BOFT is capable of producing a dense orthogonal matrix with any choice of $b$.

\textbf{Identity initialization for BOFT}. Fine-tuning approaches often commence from the pretrained model itself to ensure that subsequent updates remain close to the original parameters. For instance, LoRA applies zero initialization to its low-rank matrices. In the case of BOFT, each butterfly module is initialized as an identity matrix (\emph{i.e.}, the skew-symmetric matrix involved in the Cayley transform is set to zero at initialization).

\textbf{Multiplicative Dropout for BOFT}. While LoRA~\cite{hulora2022} incorporates a Dropout mechanism within its low-rank weight update to mitigate overfitting, and conventional Dropout~\cite{srivastava2014dropout} is inherently compatible with LoRA, such straightforward application is unsuitable for BOFT because of the multiplicative nature of its weight updates. To overcome this limitation, we introduce a variant—multiplicative Dropout—for use with BOFT. The orthogonal matrix $\bm{R}(m,b)$ comprises $m$ orthogonal butterfly components, which can be conveniently reorganized into $2b\times 2b$ block-diagonal orthogonal matrices. Our multiplicative Dropout strategy involves the random selection of $p_1$ percent of the butterfly components along with $p_2$ percent of the diagonal blocks within each component, subsequently substituting these selected blocks with identity matrices.

\section{Intriguing Insights and Discussions}

\textbf{Expressivity of BOFT}. The butterfly configuration combined with permutations is capable of exactly representing numerous classical fast linear transforms~\cite{dao2019learning,dao2019kaleidoscope} (\emph{e.g.}, fast Fourier transform, Hadamard transform). However, the extent to which our orthogonal butterfly matrix can approximate an arbitrary orthogonal matrix remains to be fully characterized. To investigate this, we perform a numerical experiment where a random dense orthogonal matrix~\cite{anderson1987generation} of dimension $512\times 512$ is approximated. The findings, presented in Figure~\ref{fig:para_eff}, reflect averages over 10 distinct random initializations. In these plots, the y-axis indicates the approximation error, while the x-axis shows the effective count of trainable parameters. Each curve in the same color corresponds to BOFT utilizing an identical block size, with the leftmost marker on each curve representing the performance of BOFT($1$,$b$) (that is, the standard block-diagonal OFT for block size $b$). Overall, BOFT tends to offer superior parameter efficiency in comparison with OFT. For instance, BOFT($9$,$2$) surpasses BOFT($1$,$16$) in expressiveness despite utilizing fewer parameters. Moreover, employing a smaller $b$ and a larger $m$ typically enhances parameter efficiency for BOFT; for example, BOFT($6$,$4$) requires substantially fewer parameters while achieving an approximation error on par with that of BOFT($2$,$16$). In summary, the butterfly structure defines a more structured and restricted subset of the orthogonal group relative to the block-diagonal framework, which underpins the provable greater expressivity of BOFT over OFT when both utilize the same block size.

\begin{theorem}[Expressivity of BOFT]\label{thm:exp_boft}
    BOFT possesses greater expressiveness than OFT when both utilize the same block size. To enable the butterfly matrix to represent any orthogonal matrix of dimension $d$, one can compose butterfly matrices as follows: $\bm{B}_{d-1,1}(d)\bm{B}^\top_{d-1,2}(d)\cdots\bm{B}_{1,1}(d)\bm{B}^\top_{1,2}(d)$, where each $\bm{B}_{i,j}(d)$ denotes a butterfly matrix for every $i$ and $j$.
\end{theorem}

Theorem~\ref{thm:exp_boft} motivates a straightforward extension of BOFT by generalizing the final orthogonal matrix to $\thickmuskip=2mu \medmuskip=2mu \bm{R}^G(m_1,b_1,m_2,b_2,l)=\bm{R}_{l,1}(m_1,b_1)\bm{R}_{l,2}^T(m_2,b_2)\cdots\bm{R}_{1,1}(m_1,b_1)\bm{R}_{1,2}^T(m_2,b_2)$, where $\bm{R}_{i,j}^T(m,b)$ indicates the specific orthogonal matrices employed within BOFT. When the parameters are set to $m_1=m_2=\log d$, $b_1=b_2=2$, and $l=d-1$, the resulting matrix $\bm{R}^G(m_1,b_1,m_2,b_2,l)$ becomes capable of representing the complete orthogonal group. This hierarchical structure is known as the kaleidoscope hierarchy~\cite{dao2019kaleidoscope}. Nonetheless, it is important to point out that enhanced representational capacity does not automatically correspond to improved performance during finetuning; for example, even though full finetuning possesses universal expressivity, it frequently fails to deliver optimal outcomes. Achieving a balance between expressivity and regularization is crucial for the generalization capacity in model finetuning. BOFT expands the search space for finetuning parameters by incorporating structural priors, thus allowing a more favorable balance between model expressiveness and regularity.

\textbf{Spectral properties}. In terms of spectral characteristics, orthogonal finetuning tends to outperform LoRA because it exactly maintains the spectral norm of the original pretrained weight matrix $\bm{W}^0$. To illustrate, recall the singular value decomposition: $\bm{W}^0=\bm{U}\bm{\Sigma}\bm{V}^\top$, where $\bm{U}$ and $\bm{V}$ are orthogonal matrices, and $\bm{\Sigma}$ contains the singular values on its diagonal. With both OFT and BOFT, an orthogonal matrix $\bm{R}$ is left-multiplied with the decomposition, resulting in adjusted weights of the form $\bm{R}\bm{U}\bm{\Sigma}\bm{V}^\top$. This operation leaves the largest singular value—i.e., the spectral norm—of $\bm{W}^0$ unchanged. Prior work has demonstrated that maintaining the spectral norm is highly advantageous for both the stability of training and model generalization~\cite{miyato2018spectral,yoshida2017spectral}. Additional noteworthy mathematical aspects are provided in Appendix~\ref{app:sn_property}.

\textbf{Orthogonal finetuning interpreted as bilinear similarity learning}. The BOFT approach can be formulated as optimizing the bilinear similarity $\bm{w}^0_i\bm{R}\bm{x}$, with $\bm{w}^0_i$ representing the $i$-th column vector (neuron) from the original weight matrix $\bm{W}^0$. From this perspective, BOFT centers around learning a bilinear similarity matrix $\bm{R}$, subject to the stringent constraint that $\bm{R}$ must be orthogonal. This perspective establishes a natural link to both distance metric learning~\cite{xing2002distance} and the theory of bilinear forms~\cite{roman2005advanced}.

\textbf{Inductive bias and generalization in BOFT}. In BOFT, $\bm{R}(m,b)$ typically denotes a structured subset of the orthogonal group, thereby imposing constraints on the class of hypotheses. This inherent structure gives rise to an inductive bias within the framework. We propose that the particular inductive bias emerging from butterfly factorization supports improved generalization, as it encapsulates the recurrent structured motifs observed in numerous well-known linear transformations~\cite{dao2019kaleidoscope}, including the discrete Fourier transform, discrete sine/cosine transforms, and the Hadamard transform. In addition, the use of sparse matrix factorization within BOFT can introduce further implicit inductive biases~\cite{gunasekar2017implicit,li2018algorithmic,liu2019neural}.

\textbf{Comparison to butterfly-based sparse training}. Numerous prior studies~\cite{dao2019learning,dao2019kaleidoscope,chen2022pixelated,dao2022monarch} have explored the use of butterfly parameterization in the context of sparse training. These works generally involve direct reparameterization of neural network weight matrices using the butterfly structure and proceed to train the models from random initialization. In contrast, \cite{dao2022monarch} explores a finetuning approach that involves first projecting pretrained weights onto a modified version of butterfly matrices and subsequently optimizing these projected parts for downstream applications. Unlike these strategies, BOFT introduces a distinct finetuning method that applies transformations to the weights via layer-shared weight matrices.

\section{Concluding Remarks and Limitations}

In this work, we introduce Orthogonal Butterfly, a broadly applicable and parameter-efficient finetuning approach for foundation models that leverages the butterfly structure. Our primary observation is that enhanced parameter efficiency can be attained by expressing a dense orthogonal matrix as the product of several sparse orthogonal matrices. To systematically discover suitable matrix decompositions, we introduce a graphical information transmission framework. Utilizing this framework, we observe that the butterfly structure aligns well with the goals of sparse orthogonal matrix decomposition. Through experiments, we provide evidence of BOFT’s empirical advantages in finetuning large language models, large vision models, and text-to-image generative models. Furthermore, our results confirm BOFT’s effectiveness as a universal method for model finetuning.

Although BOFT demonstrates strong empirical performance, it is not without limitations. The resulting orthogonal matrix in BOFT arises from the multiplication of several orthogonal matrices, which introduces slightly greater training time overhead compared to OFT. Optimizing the training efficiency of BOFT is still an unresolved issue. On the positive side, after finetuning, the orthogonal matrices learned through BOFT can be directly incorporated into the pretrained model, incurring no extra inference time. Additionally, it is still unclear if the butterfly network offers the optimal approach for information transmission. Our framework for information transmission opens up the opportunity to borrow insights from computer networking—a field where the study of efficient network topologies for information transfer is extensive. We anticipate that exploring alternative, potentially more effective network structures could lead to improved methods for constructing dense orthogonal matrices.

\end{document}